\documentclass[12pt]{article}

\usepackage{url,verbatim, amsmath, amssymb, amsfonts, dsfont, color}
\textwidth 15true cm
\textheight 9true in
\oddsidemargin 0.30true in
\evensidemargin 0.30true in
\topmargin -0.3true in
\headsep 0.4true in

\def\sgn{{\rm sign}}
\def\drightarrow{{\stackrel{d}{\rightarrow}}}
\def\dotsim{\stackrel{.}{\sim}}
\def\goesd{\stackrel{d}{\rightarrow}}
\def\goesp{\stackrel{p}{\rightarrow}}
\def\goeslaw{\stackrel{{\cal L}}{\longrightarrow}}
%
\def\vp{{\varphi}}
\def\E{{\rm E}}
\def\Pvalue{{$p$-value}}
\def\Pvalues{{$p$-values}}
\def\pr{{\rm pr}}
\def\Pr{{\rm pr}}
\def\logit{{\rm logit}}
\def\T{{ \mathrm{\scriptscriptstyle T} }}
\def\diag{{\rm diag}}  
\def\half{{\textstyle{1\over 2}}}
\def\ith{{$i{\rm th}$ }}
\def\var{{\rm var}}
\def\cl{{{\rm c}\ell}}
\def\checkx{{\checked\hspace{-6pt}\setminus}}
\def\grad{{\nabla}}
%====== NEW COMMANDS ======
\newcommand{\R}{\mathbb{R}}
\newcommand{\N}{\mathbb N}
\newcommand{\Z}{\mathbb Z}
\newcommand{\Q}{\mathbb Q}
\newcommand{\1}{\mathds 1}
\newcommand{\Var}{\text{Var}}
\newcommand{\Cov}{\text{Cov}}
\newcommand{\sd}{\text{sd}}
\newcommand{\se}{\text{se}}
\renewcommand{\L}{\mathcal {L}}
\renewcommand{\lambda}{\uplambda}
\newcommand{\iidsim}{\stackrel{iid}{\sim}}
\newcommand{\otherwise}{\text{otherwise}}
\newcommand{\indep}{\perp \!\!\! \perp}
\renewcommand{\epsilon}{\varepsilon}
\newcommand{\st}{\text{ s.t. }}
\newcommand{\ds}{\displaystyle}
\newcommand{\ind}{\stackrel{d}{\to}}
\newcommand{\inp}{\stackrel{p}{\to}}
\newcommand{\inas}{\stackrel{a.s.}{\to}}
\DeclareMathOperator*{\argmax}{arg\,max}
\DeclareMathOperator*{\argmin}{arg\,min}

\sloppy %permits line-breaking of urls

\def\blue{\color{blue}}
\def\red{\color{red}}
\def\sgn{{\rm sign}}
\def\drightarrow{{\stackrel{d}{\rightarrow}}}
\begin{document}

\noindent{\bf Questions Week 2} {\hfill STA 2212S 2026}\\

%{\bf Due January 14} \hfill{\it MS, Exercise 5.2}

\bigskip

\begin{enumerate}

\item {\em From SM Exercise 11.5.2:} 
Consider $X \sim Binom(m,\theta)$  and assume a $Beta(\alpha,\beta)$ prior for $\theta$. As shown in class, the posterior distribution $\pi(\theta\mid x)$ is $Beta(\alpha+x, \beta+m-x)$. 
\begin{enumerate}
\item Show that the marginal distribution of $X$ is 
\begin{eqnarray*}
	f(x\mid \mu, \nu) = \frac{\Gamma(\nu)}{\Gamma(\nu\mu)\Gamma\{\nu(1-\mu)\}}{m\choose x} \frac{\Gamma(x+\nu\mu)\Gamma\{m-x+\nu(1-\mu)\}}{\Gamma(m+\nu)}, \\ \quad x=0, \dots, m,
\end{eqnarray*}
where $\mu=\alpha/(\alpha+\beta)$ and $\nu=\alpha+\beta$, and deduce that
\begin{eqnarray*}
	\E(X/m) = \mu, \quad \text{var}(X/m) = \frac{\mu(1-\mu)}{m}\left(1+\frac{m-1}{\nu+1}\right).
\end{eqnarray*}
\item Suppose now that we have $n$ independent binomial responses $X_1, \dots, X_n$, each with probability of success $\theta_i$:  $$f(x_i\mid \theta_i) = {m\choose x_i}\theta_i^{x_i}(1-\theta_i)^{m-x_i}, \quad x_i = 0, \dots, m.$$ 
%\item {(c)}	
%Using the $Beta(\alpha,\beta)$ prior for $\theta_i$, show that the expected value of the posterior distribution of $\theta_i$ is \begin{equation}\E(\theta_i\mid\bs{x}) = \frac{x_i+{\alpha}}{m+{\alpha + \beta}}.\end{equation}
Describe how the results of (a) can be used to estimate the parameters $(\alpha,\beta)$ of the prior. 
\item Using these estimates, show that an estimate of the posterior expected value of $\theta_i$ is a weighted combination of the proportion of successes in the $i$th group and the proportion of successes in all groups combined. 
\end{enumerate}
\textcolor{red}{(a) The prior density of $\theta$ is 
\begin{equation*}
    \pi(\theta) = \frac{\Gamma(\alpha + \beta)}{\Gamma(\alpha)\Gamma(\beta)}\theta^{\alpha-1}(1-\theta)^{\beta - 1}
\end{equation*}
for $\theta \in \Theta = (0, 1)$. 
Then the marginal density of $X$ is 
\small
\begin{align*}
    p(x) &= \int_0^1 \mathcal L(x ; \theta)\pi(\theta) d\theta\\
    &= \int_0^1 \frac{\Gamma(\alpha+\beta)}{\Gamma(\alpha)\Gamma(\beta)} \theta^{\alpha-1}(1-\theta)^{\beta-1} \binom mx \theta^x (1-\theta)^{m-x} d\theta\\
    &= \frac{\Gamma(\alpha+\beta)}{\Gamma(\alpha)\Gamma(\beta)}\binom mx \frac{\Gamma(\alpha+x)\Gamma(\beta+m-x)}{\Gamma(\alpha+\beta+m)} \int_{\Theta}\underbrace{\frac{\Gamma(\alpha+\beta+m)}{\Gamma(\alpha+x)\Gamma(\beta+m-x)} \theta^{\alpha+x-1}(1-\theta)^{\beta+m-x-1}}_{\text{density of Beta}(\alpha+x, \beta+m+x)}d\theta\\
    &= \frac{\Gamma(\alpha+\beta)}{\Gamma(\alpha)\Gamma(\beta)} \binom mx \frac{\Gamma(\alpha+x)\Gamma(\beta+m-x)}{\Gamma(\alpha+\beta+m)}\\
    &= \frac{\Gamma(\nu)}{\Gamma(\nu\mu)\Gamma(\nu(1-\mu))}\binom mx \frac{\Gamma(x+\nu\mu)\Gamma(m-x + \nu(1-\mu))}{\Gamma(m + \nu)}
\end{align*}
\normalsize
By the Law of Total Expectation, 
\begin{equation*}
    E(X) = E_\theta[E(X \mid \theta)] = E_\theta[m\theta] = mE_\theta(\theta) = m \mu
\end{equation*}
thus $E(X / m) = \mu$. By the Law of Total Variance, 
\begin{align*}
    \Var(X) &= E_\theta[\Var(X \mid \theta)] + \Var_\theta[E(X \mid \theta)]\\
    &= E_\theta[m\theta(1-\theta)] + \Var_\theta[m\theta]\\
    &= mE_\theta[\theta - \theta^2] + m^2\frac{\mu(1-\mu)}{\nu + 1}\\
    &= m(\mu - \Var_\theta(\theta) + \mu^2) + m^2\frac{\mu(1-\mu)}{\nu + 1}\\
    &= m(\mu(1-\mu) - \Var(\theta)) + m^2\frac{\mu(1-\mu)}{\nu + 1}\\
    &= m\mu(1-\mu) + (m^2 - m) \frac{\mu(1-\mu)}{\nu + 1}
\end{align*}
so 
\begin{equation*}
    \Var(X/m) = \frac{1}{m^2} \Var(X) = \frac{\mu(1-\mu)}{m}\left( 1 + \frac{m-1}{\nu + 1}\right)
\end{equation*}
(b) The result from (a) implies that we have 
\begin{equation*}
    f(x_i \mid \alpha, \beta) = \frac{\Gamma(\nu)}{\Gamma(\nu\mu)\Gamma(\nu(1-\mu))}\binom{m}{x_i} \frac{\Gamma(x_i+\nu\mu)\Gamma(m-x_i + \nu(1-\mu))}{\Gamma(m + \nu)}
\end{equation*}
Define the likelihood function 
\begin{equation*}
    \mathcal L(\alpha, \beta \mid \mathbf x) = \prod_{i=1}^n \frac{\Gamma(\nu)}{\Gamma(\nu\mu)\Gamma(\nu(1-\mu))}\binom{m}{x_i} \frac{\Gamma(x_i+\nu\mu)\Gamma(m-x_i + \nu(1-\mu))}{\Gamma(m + \nu)}
\end{equation*}
where $\mathbf x = (x_1,\ldots, x_n)'$ is the vector of realizations of $\mathbf X = (X_1,\ldots, X_n)'$. 
By optimizing $\mathcal L$ and obtaining the MLEs of $\alpha$ and $\beta$ respectively, we can estimate the parameters of the prior of $\theta$. \\
(c) Using these estimates, the estimated posterior for $\theta_i$ is $\text{Beta}(\hat\alpha + x_i, \hat\beta + m -x_i)$ where $\hat\alpha$ and $\hat\beta$ are the MLEs of $\alpha$ and $\beta$ respectively. The estimated expectation is then 
\begin{equation*}
    \widehat{E(\theta_i \mid X_i)} = \frac{\hat\alpha + X_i }{(\hat\alpha + X_i) + (\hat\beta + m - X_i)} = \frac{\hat\alpha + X_i}{\hat\alpha + \hat\beta + m}
\end{equation*}
By invariance the MLE, we can estimate $\mu$ and $\nu$ as $\hat\nu = \hat\alpha + \hat\beta$ and $\hat\mu = \hat\alpha/(\hat\alpha + \hat\beta)$, thus 
\begin{equation*}
    \widehat{E(\theta_i \mid X_i)} = \frac{m}{m + \hat\nu}\frac{X_i}{m} + \frac{\hat\nu}{m + \hat\nu}\hat\mu
\end{equation*}
which is a weighted sum of the proportion of successes in the $i$'th group ($X_i / m$), and the overall proportion $\hat\mu$. 
}

\item {\it Cox \& Hinkley, Ex.8.1}: This example illustrates in extreme form why exact unbiasedness of an estimator is not always useful. Suppose $X$ follows the geometric distribution with probability mass function 
$$
f(x;\theta) = (1-\theta)^x \theta, \quad 0 \le\theta\le 1,\, x = 0, 1, \dots \,.
$$

Show that the unique unbiased estimator of $\theta$ is 
$\widetilde\theta (x) = 1$ if $x=0$, and $\widetilde\theta(x) = 0$ for $x = 1, 2, \dots$\,.  

\textcolor{red}{
    First, we show that $\tilde\theta$ is unbiased. 
    \begin{equation*}
        E(\tilde\theta) = \sum_{x=0}^\infty \tilde\theta(x)(1-\theta)^x\theta = (1-\theta)^0\theta = \theta
    \end{equation*}
    hence $\tilde \theta$ is unbiased. 
    To show uniqueness, suppose there exists another estimator $\hat\theta$ such that $\hat\theta$ is unbiased. This implies 
    \begin{equation*}
        \theta = E(\hat\theta) = \sum_{x=0}^\infty \hat\theta(x) (1-\theta)^x \theta = \hat\theta(0)\theta + \sum_{x=1}^\infty \hat\theta(x)(1-\theta)^x\theta
    \end{equation*}
    thus 
    \begin{equation*}
        1 = \hat\theta(0) + \sum_{x=1}^\infty \hat\theta(x) (1-\theta)^x
    \end{equation*}
    Since for each $x \geq 1$, $\hat\theta(x)(1-\theta)^x$ is a term dependent on $\theta$ unless $\hat\theta(x) = 0$, then we must have $\hat\theta(x) = 0$ for all $x = 1,\ldots,$, which implies $\hat\theta(0) = 1$. This implies $\hat\theta = \tilde\theta$, hence showing uniqueness. 
}
\end{enumerate}
\end{document}

